% block: references-heading page:223 type:chapter
\clearpage\section*{{参考文献}}\phantomsection
\makeatletter\def\@currentlabel{References}\label{chap:references}\makeatother
\begin{enumerate}[label=\arabic*.,ref=\arabic*,leftmargin=3em]

% block: reference-source-001 page:223 type:ordinary-paragraph
\item\label{bib:source-1}Abe, K., Fujino, S.: Converting BiCR method for linear equations with complex symmetric matrices. App. Math. Comput. 321, 564–576 (2018)

% block: reference-source-002 page:223 type:ordinary-paragraph
\item\label{bib:source-2}Abe, K., Sleijpen, G.L.G.: BiCR variants of the hybrid BiCG methods for solving linear systems with nonsymmetric matrices. J. Comput. Appl. Math. 234, 985–994 (2010)

% block: reference-source-003 page:223 type:ordinary-paragraph
\item\label{bib:source-3}Absil, P.-A., Mahony, R., Sepulchre, R.: Optimization Algorithms on Matrix Manifolds. Princeton University Press, Princeton (2008)

% block: reference-source-004 page:223 type:ordinary-paragraph
\item\label{bib:source-4}Aihara, K.: GPBi-CGstab(L): A Lanczos-type product method unifying Bi-CGstab(L) and GPBi-CG. Numer. Linear Algebra Appl. 27, e2298 (2020)

% block: reference-source-005 page:223 type:ordinary-paragraph
\item\label{bib:source-5}Aihara, K., Abe, K., Ishiwata, E.: An alternative implementation of the IDRstab method saving vector updates. JSIAM Lett. 3, 69–72 (2011)

% block: reference-source-006 page:223 type:ordinary-paragraph
\item\label{bib:source-6}Aihara, K., Abe, K., Ishiwata, E.: A quasi-minimal residual variant of IDRstab using the residual smoothing technique. Appl. Math. Comput. 236, 67–77 (2014)

% block: reference-source-007 page:223 type:ordinary-paragraph
\item\label{bib:source-7}Aihara, K., Abe, K., Ishiwata, E.: A variant of IDRstab with reliable update strategies for solving sparse linear systems. J. Comput. Appl. Math. 259, 244–258 (2014)

% block: reference-source-008 page:223 type:ordinary-paragraph
\item\label{bib:source-8}Al-Mohy, A.H., Higham, N.J.: A new scaling and squaring algorithm for the matrix exponential. SIAM J. Matrix Anal. Appl. 31, 970–989 (2009)

% block: reference-source-009 page:223 type:ordinary-paragraph
\item\label{bib:source-9}Al-Mohy, A.H., Higham, N.J.: Improved inverse scaling and squaring algorithms for the matrix logarithm. SIAM J. Sci. Comput. 34, C153–C169 (2012)

% block: reference-source-010 page:223 type:ordinary-paragraph
\item\label{bib:source-10}Axelsson, O.: Iterative Solution Methods. Cambridge University Press, Cambridge (1994)

% block: reference-source-011 page:223 type:ordinary-paragraph
\item\label{bib:source-11}Baglama, J., Reichel, L.: Augmented GMRES-type methods. Numer. Linear Algebra Appl. 14, 337–350 (2007)

% block: reference-source-012 page:223 type:ordinary-paragraph
\item\label{bib:source-12}Bai, Z., Demmel, J., Dongarra, J., Ruhe, A., van der Vorst, H.A. (eds.): Templates for the Solution of Algebraic Eigenvalue Problems. SIAM, Philadelphia (2000)

% block: reference-source-013 page:223 type:ordinary-paragraph
\item\label{bib:source-13}Baker, A.H., Dennis, J.M., Jessup, E.R.: On improving linear solver performance: a block variant of GMRES. SIAM J. Sci. Comput. 27, 1608–1626 (2006)

% block: reference-source-014 page:223 type:ordinary-paragraph
\item\label{bib:source-14}Baker, A.H., Jessup, E.R., Kolev, T.V.: A simple strategy for varying the restart parameter in GMRES(m). J. Comput. Appl. Math. 230, 751–761 (2009)

% block: reference-source-015 page:223 type:ordinary-paragraph
\item\label{bib:source-15}Baker, A.H., Jessup, E.R., Manteuffel, T.: A technique for accelerating the convergence of restarted GMRES. SIAM J. Matrix Anal. Appl. 26, 962–984 (2005)

% block: reference-source-016 page:223 type:ordinary-paragraph
\item\label{bib:source-16}Bank, R.E., Chan, T.F.: An analysis of the composite step biconjugate gradient method. Numer. Math. 66, 295–319 (1993)

% block: reference-source-017 page:224 type:ordinary-paragraph
\item\label{bib:source-17}Bank, R.E., Chan, T.F.: A composite step bi-conjugate gradient algorithm for nonsymmetric linear systems. Numer. Algorithms 7, 1–16 (1994)

% block: reference-source-018 page:224 type:ordinary-paragraph
\item\label{bib:source-18}Barrett, R., Berry, M., Chan, T.F., Demmel, J., Donato, J., Dongarra, J., Eijkhout, V., Pozo, R., Romine, C., van der Vorst, H.: Templates for the Solution of Linear Systems: Building Blocks for Iterative Methods. SIAM, Philadelphia (1994)

% block: reference-source-019 page:224 type:ordinary-paragraph
\item\label{bib:source-19}Bayliss, A., Goldstein, C., Turkel, E.: An iterative method for the Helmholtz equation. J. Comput. Phys. 49, 443–457 (1983)

% block: reference-source-020 page:224 type:ordinary-paragraph
\item\label{bib:source-20}Benzi, M.: Preconditioning techniques for large linear systems: a survey. J. Comput. Phys. 182, 418–477 (2002)

% block: reference-source-021 page:224 type:ordinary-paragraph
\item\label{bib:source-21}Benzi, M., Cullum, J.K., Tůma, M.: Robust approximate inverse preconditioning for the conjugate gradient method. SIAM J. Sci. Comput. 22, 1318–1332 (2000)

% block: reference-source-022 page:224 type:ordinary-paragraph
\item\label{bib:source-22}Benzi, M., Meyer, C.D., Tůma, M.: A sparse approximate inverse preconditioner for the conjugate gradient method. SIAM J. Sci. Comput. 17, 1135–1149 (1996)

% block: reference-source-023 page:224 type:ordinary-paragraph
\item\label{bib:source-23}Benzi, M., Szyld, D.B., van Duin, A.: Orderings for incomplete factorization preconditioning of nonsymmetric problems. SIAM J. Sci. Comput. 20, 1652–1670 (1999)

% block: reference-source-024 page:224 type:ordinary-paragraph
\item\label{bib:source-24}Benzi, M., Tůma, M.: A sparse approximate inverse preconditioner for nonsymmetric linear systems. SIAM J. Sci. Comput. 19, 968–994 (1998)

% block: reference-source-025 page:224 type:ordinary-paragraph
\item\label{bib:source-25}Benzi, M., Tůma, M.: A robust preconditioner with low memory requirements for large sparse least squares problems. SIAM J. Sci. Comput. 25, 499–512 (2003)

% block: reference-source-026 page:224 type:ordinary-paragraph
\item\label{bib:source-26}Boyd, S., Vandenberghe, L.: Convex Optimization. Cambridge University Press, Cambridge (2004)

% block: reference-source-027 page:224 type:ordinary-paragraph
\item\label{bib:source-27}Bramble, J.H.: Multigrid Methods. Longman Scientific and Technical, Harlow (1993)

% block: reference-source-028 page:224 type:ordinary-paragraph
\item\label{bib:source-28}Briggs, W.L., Henson, V.E., McCormick, S.F.: A Multigrid Tutorial, 2nd edn. SIAM, Philadelphia (2000)

% block: reference-source-029 page:224 type:ordinary-paragraph
\item\label{bib:source-29}Bruaset, A.M.: A Survey of Preconditioned Iterative Methods. Routledge, Boca Raton (1995)

% block: reference-source-030 page:224 type:ordinary-paragraph
\item\label{bib:source-30}Bunse-Gerstner, A., Stöver, R.: On a conjugate gradient-type method for solving complex symmetric linear systems. Linear Algebra Appl. 287, 105–123 (1999)

% block: reference-source-031 page:224 type:ordinary-paragraph
\item\label{bib:source-31}Cabral, J.C., Schaerer, C.E., Bhaya, A.: Improving GMRES(m) using an adaptive switching controller. Numer. Linear Algebra Appl. 27, e2305 (2020)

% block: reference-source-032 page:224 type:ordinary-paragraph
\item\label{bib:source-32}Carpentieri, B., Jing, Y.-F., Huang, T.-Z.: The BiCOR and CORS iterative algorithms for solving nonsymmetric linear systems. SIAM J. Sci. Comput. 33, 3020–3036 (2011)

% block: reference-source-033 page:224 type:ordinary-paragraph
\item\label{bib:source-33}Cayley, A.: A memoir on the theory of matrices. Philos. Trans. R. Soc. 148, 17–37 (1858)

% block: reference-source-034 page:224 type:ordinary-paragraph
\item\label{bib:source-34}Chan, T.F., Gallopoulos, E., Simoncini, V., Szeto, T., Tong, C.H.: A quasi-minimal residual variant of the Bi-CGSTAB algorithm for nonsymmetric systems. SIAM J. Sci. Comput. 15, 338–347 (1994)

% block: reference-source-035 page:224 type:ordinary-paragraph
\item\label{bib:source-35}Chan, T.F., Szeto, T.: Composite step product methods for solving nonsymmetric linear systems. SIAM J. Sci. Comput. 17, 1491–1508 (1996)

% block: reference-source-036 page:224 type:ordinary-paragraph
\item\label{bib:source-36}Chiba, K.: Gyouretsu no Kansu to Jordan Hyoujun Kei (in Japanese). Scientist Press, Tokyo (2010)

% block: reference-source-037 page:224 type:ordinary-paragraph
\item\label{bib:source-37}Choi, S.-C.: Minimal residual methods for complex symmetric, skew symmetric, and skew Hermitian systems. Tech. rep., Comput. Inst., Univ. Chicago, Chicago, IL, USA (2013)

% block: reference-source-038 page:224 type:ordinary-paragraph
\item\label{bib:source-38}Chronopoulos, A.T., Ma, S.: On squaring Krylov subspace iterative methods for nonsymmetric linear systems. Tech. Rep. TR 89-67 (1989)

% block: reference-source-039 page:224 type:ordinary-paragraph
\item\label{bib:source-39}Cipra, B.A.: The best of the 20th century: Editors name top 10 algorithms. SIAM News 33, 1–2 (2000)

% block: reference-source-040 page:224 type:ordinary-paragraph
\item\label{bib:source-40}Clemens, M., Weiland, T., Rienen, U.V.: Comparison of Krylov-type methods for complex linear systems applied to high-voltage problems. IEEE Trans. Magn. 34, 3335–3338 (1998)

% block: reference-source-041 page:224 type:ordinary-paragraph
\item\label{bib:source-41}Craig, E.J.: The N -step iteration procedures. J. Math. Phys. 34, 64–73 (1955)

% block: reference-source-042 page:224 type:ordinary-paragraph
\item\label{bib:source-42}Cuthill, E., McKee, J.: Reducing the bandwidth of sparse symmetric matrices. In: Proceedings of the 1969 24th National Conference, pp. 157–172. Association for Computing Machinery, New York, NY, USA (1969)

% block: reference-source-043 page:224 type:ordinary-paragraph
\item\label{bib:source-43}Datta, B.N., Saad, Y.: Arnoldi methods for large Sylvester-like observer matrix equations, and an associated algorithm for partial spectrum assignment. Linear Algebra Appl. 154–156, 225–244 (1991)

% block: reference-source-044 page:225 type:ordinary-paragraph
\item\label{bib:source-44}Davis, T.A.: Direct Methods for Sparse Linear Systems. SIAM, Philadelphia (2006)

% block: reference-source-045 page:225 type:ordinary-paragraph
\item\label{bib:source-45}Demmel, J.W.: Applied Numerical Linear Algebra. SIAM, Philadelphia (1997)

% block: reference-source-046 page:225 type:ordinary-paragraph
\item\label{bib:source-46}Du, L., Sogabe, T., Yu, B., Yamamoto, Y., Zhang, S.-L.: A block IDR(s) method for nonsymmetric linear systems with multiple right-hand sides. J. Comput. Appl. Math. 235, 4095–4106 (2011)

% block: reference-source-047 page:225 type:ordinary-paragraph
\item\label{bib:source-47}Du, L., Sogabe, T., Zhang, S.-L.: A variant of the IDR(s) method with the quasi-minimal residual strategy. J. Comput. Appl. Math. 236, 621–630 (2011)

% block: reference-source-048 page:225 type:ordinary-paragraph
\item\label{bib:source-48}Du, L., Sogabe, T., Zhang, S.-L.: IDR(s) for solving shifted nonsymmetric linear systems. J. Comput. Appl. Math. 274, 35–43 (2015)

% block: reference-source-049 page:225 type:ordinary-paragraph
\item\label{bib:source-49}Du, Y.-S., Hayami, K., Zheng, N., Morikuni, K., Yin, J.-F.: Kaczmarz-type inner-iteration preconditioned flexible GMRES methods for consistent linear systems. SIAM J. Sci. Comput. 43, S345–S366 (2021)

% block: reference-source-050 page:225 type:ordinary-paragraph
\item\label{bib:source-50}Dubois, P.F., Greenbaum, A., Rodrigue, G.H.: Approximating the inverse of a matrix for use in iterative algorithms on vector processors. Computing 22, 257–268 (1979)

% block: reference-source-051 page:225 type:ordinary-paragraph
\item\label{bib:source-51}Duff, I.S., Erisman, A.M., Reid, J.K.: Direct Methods for Sparse Matrices. Oxford University Press, Oxford (1986)

% block: reference-source-052 page:225 type:ordinary-paragraph
\item\label{bib:source-52}Duff, I.S., Meurant, G.A.: The effect of ordering on preconditioned conjugate gradients. BIT Numer. Math. 29, 635–657 (1989)

% block: reference-source-053 page:225 type:ordinary-paragraph
\item\label{bib:source-53}Eisenstat, S.C.: Efficient implementation of a class of preconditioned conjugate gradient methods. SIAM J. Sci. Stat. Comput. 2, 1–4 (1981)

% block: reference-source-054 page:225 type:ordinary-paragraph
\item\label{bib:source-54}Eisenstat, S.C., Elman, H.C., Schultz, M.H.: Variational iterative methods for nonsymmetric systems of linear equations. SIAM J. Numer. Anal. 20, 345–357 (1983)

% block: reference-source-055 page:225 type:ordinary-paragraph
\item\label{bib:source-55}Elbouyahyaoui, L., Heyouni, M., Tajaddini, A., Saberi-Movahed, F.: On restarted and deflated block FOM and GMRES methods for sequences of shifted linear systems. Numer. Algorithms 87, 1257–1299 (2021)

% block: reference-source-056 page:225 type:ordinary-paragraph
\item\label{bib:source-56}Eldén, L., Savas, B.: A Newton-Grassmann method for computing the best multilinear rank-$(r_1,r_2,r_3)$ approximation of a tensor. SIAM J. Matrix Anal. Appl. 31, 248–271 (2009)

% block: reference-source-057 page:225 type:ordinary-paragraph
\item\label{bib:source-57}Faber, V., Manteuffel, T.: Necessary and sufficient conditions for the existence of a conjugate gradient method. SIAM J. Numer. Anal. 21, 352–362 (1984)

% block: reference-source-058 page:225 type:ordinary-paragraph
\item\label{bib:source-58}Fletcher, R.: Conjugate gradient methods for indefinite systems. In: Watson, G.A. (ed.) Numerical Analysis, pp. 73–89. Springer, Berlin, Heidelberg (1976)

% block: reference-source-059 page:225 type:ordinary-paragraph
\item\label{bib:source-59}Forsythe, G.E., Straus, E.G.: On best conditioned matrices. Proc. Amer. Math. Soc. 6, 340–345 (1955)

% block: reference-source-060 page:225 type:ordinary-paragraph
\item\label{bib:source-60}Freund, R.W.: On conjugate gradient type methods and polynomial preconditioners for a class of complex non-Hermitian matrices. Numer. Math. 57, 285–312 (1990)

% block: reference-source-061 page:225 type:ordinary-paragraph
\item\label{bib:source-61}Freund, R.W.: Conjugate gradient-type methods for linear systems with complex symmetric coefficient matrices. SIAM J. Sci. Stat. Comput. 13, 425–448 (1992)

% block: reference-source-062 page:225 type:ordinary-paragraph
\item\label{bib:source-62}Freund, R.W.: Solution of shifted linear systems by quasi-minimal residual iterations. In: Reichel, L., Ruttan, A., Varga, R.S. (eds.) Numerical Linear Algebra and Scientific Computation, pp. 101–122. De Gruyter (1993)

% block: reference-source-063 page:225 type:ordinary-paragraph
\item\label{bib:source-63}Freund, R.W.: A transpose-free quasi-minimal residual algorithm for non-Hermitian linear systems. SIAM J. Sci. Comput. 14, 470–482 (1993)

% block: reference-source-064 page:225 type:ordinary-paragraph
\item\label{bib:source-64}Freund, R.W., Golub, G.H., Nachtigal, N.M.: Iterative solution of linear systems. Acta Numer. 1, 57–100 (1992)

% block: reference-source-065 page:225 type:ordinary-paragraph
\item\label{bib:source-65}Freund, R.W., Malhotra, M.: A block QMR algorithm for non-Hermitian linear systems with multiple right-hand sides. Linear Algebra Appl. 254, 119–157 (1997)

% block: reference-source-066 page:225 type:ordinary-paragraph
\item\label{bib:source-66}Freund, R.W., Nachtigal, N.M.: QMR: a quasi-minimal residual method\newline{}for non-Hermitian linear systems. Numer. Math. 60, 315–339 (1991)

% block: reference-source-067 page:225 type:ordinary-paragraph
\item\label{bib:source-67}Frommer, A.: BiCGStab($\ell$) for families of shifted linear systems. Computing 70, 87–109 (2003)

% block: reference-source-068 page:225 type:ordinary-paragraph
\item\label{bib:source-68}Frommer, A., Glässner, U.: Restarted GMRES for shifted linear systems. SIAM J. Sci. Comput. 19, 15–26 (1998)

% block: reference-source-069 page:226 type:ordinary-paragraph
\item\label{bib:source-69}Frommer, A., Lippert, T., Medeke, B., Schilling, K. (eds.): Numerical challenges in lattice quantum chromodynamics: joint interdisciplinary workshop of John von Neumann institute for computing, Jülich, and Institute of Applied Computer Science, Wuppertal University, August 1999. Springer-Verlag, Berlin Heidelberg (2000)

% block: reference-source-070 page:226 type:ordinary-paragraph
\item\label{bib:source-70}Frommer, A., Maass, P.: Fast CG-based methods for Tikhonov-Phillips regularization. SIAM J. Sci. Comput. 20, 1831–1850 (1999)

% block: reference-source-071 page:226 type:ordinary-paragraph
\item\label{bib:source-71}Fujino, S.: GPBiCG(m, l): A hybrid of BiCGSTAB and GPBiCG methods with efficiency and robustness. Appl. Numer. Math. 41, 107–117 (2002)

% block: reference-source-072 page:226 type:ordinary-paragraph
\item\label{bib:source-72}Fujino, S., Zhang, S.-L.: Hanpukuhou no Suri (in Japanese). Asakura Shoten, Tokyo (1996)

% block: reference-source-073 page:226 type:ordinary-paragraph
\item\label{bib:source-73}Fujiwara, T., Hoshi, T., Yamamoto, S., Sogabe, T., Zhang, S.-L.: Novel algorithm of large-scale simultaneous linear equations. J. Phys.: Condens. Matter 22, 074206 (2010)

% block: reference-source-074 page:226 type:ordinary-paragraph
\item\label{bib:source-74}Gaul, A., Gutknecht, M.H., Liesen, J., Nabben, R.: A framework for deflated and augmented Krylov subspace methods. SIAM J. Matrix Anal. Appl. 34, 495–518 (2013)

% block: reference-source-075 page:226 type:ordinary-paragraph
\item\label{bib:source-75}George, A.: Computer implementation of the finite element method. Tech. Rep. STAN-CS-208, Stanford University, Stanford, CA, USA (1971)

% block: reference-source-076 page:226 type:ordinary-paragraph
\item\label{bib:source-76}George, A.: Nested dissection of a regular finite element mesh. SIAM J. Numer. Anal. 10, 345–363 (1973)

% block: reference-source-077 page:226 type:ordinary-paragraph
\item\label{bib:source-77}George, A., Liu, J.W.: The evolution of the minimum degree ordering algorithm. SIAM Rev. 31, 1–19 (1989)

% block: reference-source-078 page:226 type:ordinary-paragraph
\item\label{bib:source-78}Golub, G., Kahan, W.: Calculating the singular values and pseudo-inverse of a matrix. J. Soc. Ind. Appl. Math. Ser. B Numer. Anal. 2, 205–224 (1965)

% block: reference-source-079 page:226 type:ordinary-paragraph
\item\label{bib:source-79}Golub, G.H., Loan, C.F.V.: Matrix Computations, 4th edn. Johns Hopkins University Press, Baltimore (2013)

% block: reference-source-080 page:226 type:ordinary-paragraph
\item\label{bib:source-80}Golub, G.H., O’Leary, D.P.: Some history of the conjugate gradient and Lanczos algorithms: 1948–1976. SIAM Rev. 31, 50–102 (1989)

% block: reference-source-081 page:226 type:ordinary-paragraph
\item\label{bib:source-81}Greenbaum, A.: Iterative Methods for Solving Linear Systems. SIAM, Philadelphia (1997)

% block: reference-source-082 page:226 type:ordinary-paragraph
\item\label{bib:source-82}Grote, M.J., Huckle, T.: Parallel preconditioning with sparse approximate inverses. SIAM J. Sci. Comput. 18, 838–853 (1997)

% block: reference-source-083 page:226 type:ordinary-paragraph
\item\label{bib:source-83}Gu, X.-M., Clemens, M., Huang, T.-Z., Li, L.: The SCBiCG class of algorithms for complex symmetric linear systems with applications in several electromagnetic model problems. Comput. Phys. Commun. 191, 52–64 (2015)

% block: reference-source-084 page:226 type:ordinary-paragraph
\item\label{bib:source-84}Gu, X.-M., Huang, T.-Z., Li, L., Li, H., Sogabe, T., Clemens, M.: Quasi-minimal residual variants of the COCG and COCR methods for complex symmetric linear systems in electromagnetic simulations. IEEE Trans. Microw. Theory Tech. 62, 2859–2867 (2014)

% block: reference-source-085 page:226 type:ordinary-paragraph
\item\label{bib:source-85}Gu, X.-M., Huang, T.-Z., Meng, J., Sogabe, T., Li, H.B., Li, L.: BiCR-type methods for families of shifted linear systems. Comput. Math. Appl. 68, 746–758 (2014)

% block: reference-source-086 page:226 type:ordinary-paragraph
\item\label{bib:source-86}Guennouni, A.E., Jbilou, K., Sadok, H.: A block version of BiCGSTAB for linear systems with multiple right-hand sides. Electron. Trans. Numer. Anal. 16, 129–142 (2003)

% block: reference-source-087 page:226 type:ordinary-paragraph
\item\label{bib:source-87}Gutknecht, M.H.: Variants of BiCGSTAB for matrices with complex spectrum. SIAM J. Sci. Comput. 14, 1020–1033 (1993)

% block: reference-source-088 page:226 type:ordinary-paragraph
\item\label{bib:source-88}Gutknecht, M.H.: IDR explained. Electron. Trans. Numer. Anal. 36, 126–148 (2009)

% block: reference-source-089 page:226 type:ordinary-paragraph
\item\label{bib:source-89}Gutknecht, M.H.: Spectral deflation in Krylov solvers: a theory of coordinate space based methods. Electron. Trans. Numer. Anal. 39, 156–185 (2012)

% block: reference-source-090 page:226 type:ordinary-paragraph
\item\label{bib:source-90}Gutknecht, M.H.: Deflated and augmented Krylov subspace methods: a framework for deflated BiCG and related solvers. SIAM J. Matrix Anal. Appl. 35, 1444–1466 (2014)

% block: reference-source-091 page:226 type:ordinary-paragraph
\item\label{bib:source-91}Hackbusch, W.: Multi-Grid Methods and Applications. Springer-Verlag, Berlin Heidelberg (1985)

% block: reference-source-092 page:226 type:ordinary-paragraph
\item\label{bib:source-92}Hackbusch, W.: Iterative Solution of Large Sparse Systems of Equations, Applied Mathematical Sciences, vol. 95. Springer-Verlag, New York (1994)

% block: reference-source-093 page:226 type:ordinary-paragraph
\item\label{bib:source-93}Hadjidimos, A.: Successive overrelaxation (SOR) and related methods. J. Comput. Appl. Math. 123, 177–199 (2000)

% block: reference-source-094 page:226 type:ordinary-paragraph
\item\label{bib:source-94}Hayami, K., Yin, J.-F., Ito, T.: GMRES methods for least squares problems. SIAM J. Matrix Anal. Appl. 31, 2400–2430 (2010)

% block: reference-source-095 page:227 type:ordinary-paragraph
\item\label{bib:source-95}Hestenes, M.R., Stiefel, E.: Methods of conjugate gradients for solving linear systems. J. Res. Natl. Bur. Stand. 49, 409–436 (1952)

% block: reference-source-096 page:227 type:ordinary-paragraph
\item\label{bib:source-96}Higham, N.J.: Newton’s method for the matrix square root. Math. Comp. 46, 537–549 (1986)

% block: reference-source-097 page:227 type:ordinary-paragraph
\item\label{bib:source-97}Higham, N.J.: Functions of Matrices. SIAM, Philadelphia (2008)

% block: reference-source-098 page:227 type:ordinary-paragraph
\item\label{bib:source-98}Higham, N.J.: The scaling and squaring method for the matrix exponential revisited. SIAM Rev. 51, 747–764 (2009)

% block: reference-source-099 page:227 type:ordinary-paragraph
\item\label{bib:source-99}Hochbruck, M., Lubich, C.: On Krylov subspace approximations to the matrix exponential operator. SIAM J. Numer. Anal. 34, 1911–1925 (1997)

% block: reference-source-100 page:227 type:ordinary-paragraph
\item\label{bib:source-100}Hoshi, T., Kawamura, M., Yoshimi, K., Motoyama, Y., Misawa, T., Yamaji, Y., Todo, S., Kawashima, N., Sogabe, T.: K$\omega$ - Open-source library for the shifted Krylov subspace method of the form $(zI-H)\boldsymbol{x}=\boldsymbol{b}$. Comput. Phys. Commun. 258, 107536 (2021)

% block: reference-source-101 page:227 type:ordinary-paragraph
\item\label{bib:source-101}Iannazzo, B.: A note on computing the matrix square root. Calcolo 40, 273–283 (2003)

% block: reference-source-102 page:227 type:ordinary-paragraph
\item\label{bib:source-102}Iannazzo, B.: On the Newton method for the matrix pth root. SIAM J. Matrix Anal. Appl. 28, 503–523 (2006)

% block: reference-source-103 page:227 type:ordinary-paragraph
\item\label{bib:source-103}Imakura, A., Li, R.-C., Zhang, S.-L.: Locally optimal and heavy ball GMRES methods. Japan J. Indust. Appl. Math. 33, 471–499 (2016)

% block: reference-source-104 page:227 type:ordinary-paragraph
\item\label{bib:source-104}Imakura, A., Sogabe, T., Zhang, S.-L.: An efficient variant of the GMRES(m) method based on the error equations. E. Asian J. App. Math. 2, 19–32 (2012)

% block: reference-source-105 page:227 type:ordinary-paragraph
\item\label{bib:source-105}Imakura, A., Sogabe, T., Zhang, S.-L.: An efficient variant of the restarted shifted GMRES method for solving shifted linear systems. J. Math. Res. Appl. 33, 127–141 (2013)

% block: reference-source-106 page:227 type:ordinary-paragraph
\item\label{bib:source-106}Imakura, A., Sogabe, T., Zhang, S.-L.: A look-back-type restart for the restarted Krylov subspace methods for solving non-Hermitian linear systems. Japan J. Indust. Appl. Math. 35, 835–859 (2018)

% block: reference-source-107 page:227 type:ordinary-paragraph
\item\label{bib:source-107}Ipsen, I.C.F.: Numerical Matrix Analysis: Linear Systems and Least Squares. SIAM, Philadelphia (2009)

% block: reference-source-108 page:227 type:ordinary-paragraph
\item\label{bib:source-108}Ishikawa, K.-I., Sogabe, T.: A thick-restart Lanczos type method for Hermitian J-symmetric eigenvalue problems. Japan J. Indust. Appl. Math. 38, 233–256 (2021)

% block: reference-source-109 page:227 type:ordinary-paragraph
\item\label{bib:source-109}Itoh, S.: Improvement of preconditioned bi-Lanczos-type algorithms with residual norm minimization for the stable solution of systems of linear equations. Japan J. Indust. Appl. Math. 39, 17–74 (2022)

% block: reference-source-110 page:227 type:ordinary-paragraph
\item\label{bib:source-110}Iwase, S., Hoshi, T., Ono, T.: Numerical solver for first-principles transport calculation based on real-space finite-difference method. Phys. Rev. E 91, 063305 (2015)

% block: reference-source-111 page:227 type:ordinary-paragraph
\item\label{bib:source-111}Jin, S., Bulgac, A., Roche, K., Wlazłowski, G.: Coordinate-space solver for superfluid many-fermion systems with the shifted conjugate-orthogonal conjugate-gradient method. Phys. Rev. C 95, 044302 (2017)

% block: reference-source-112 page:227 type:ordinary-paragraph
\item\label{bib:source-112}Jing, Y., Huang, T.: Restarted weighted full orthogonalization method for shifted linear systems. Comput. Math. Appl. 57, 1583–1591 (2009)

% block: reference-source-113 page:227 type:ordinary-paragraph
\item\label{bib:source-113}Jing, Y., Huang, T., Carpentieri, B., Duan, Y.: Exploiting the composite step strategy to the biconjugate A-orthogonal residual method for non-Hermitian linear systems. J. Appl. Math. 2013, 16pp. (2013)

% block: reference-source-114 page:227 type:ordinary-paragraph
\item\label{bib:source-114}Jing, Y., Huang, T.-Z., Zhang, Y., Li, L., Cheng, G.-H., Ren, Z.-G., Duan, Y., Sogabe, T., Carpentieri, B.: Lanczos-type variants of the COCR method for complex nonsymmetric linear systems. J. Comput. Phys. 228, 6376–6394 (2009)

% block: reference-source-115 page:227 type:ordinary-paragraph
\item\label{bib:source-115}Jing, Y., Yuan, P., Huang, T.: A simpler GMRES and its adaptive variant for shifted linear systems. Numer. Linear Algebra Appl. 24, e2076 (2017)

% block: reference-source-116 page:227 type:ordinary-paragraph
\item\label{bib:source-116}Johnson, O.G., Micchelli, C.A., Paul, G.: Polynomial preconditioners for conjugate gradient calculations. SIAM J. Numer. Anal. 20, 362–376 (1983)

% block: reference-source-117 page:227 type:ordinary-paragraph
\item\label{bib:source-117}Joubert, W.: On the convergence behavior of the restarted GMRES algorithm for solving nonsymmetric linear systems. Numer. Linear Algebra Appl. 1, 427–447 (1994)

% block: reference-source-118 page:227 type:ordinary-paragraph
\item\label{bib:source-118}Kikuchi, F.: Yugen Yosoho Gaisetsu (in Japanese). Saiensu-Sha, Tokyo (1980)

% block: reference-source-119 page:227 type:ordinary-paragraph
\item\label{bib:source-119}Lanczos, C.: Solution of systems of linear equations by minimized iterations. J. Res. Natl. Bur. Stand. 49, 33 (1952)

% block: reference-source-120 page:227 type:ordinary-paragraph
\item\label{bib:source-120}Lee, D., Hoshi, T., Sogabe, T., Miyatake, Y., Zhang, S.-L.: Solution of the k-th eigenvalue problem in large-scale electronic structure calculations. J. Comput. Phys. 371, 618–632 (2018)

% block: reference-source-121 page:228 type:ordinary-paragraph
\item\label{bib:source-121}Liesen, J., Strakoš, Z.: On optimal short recurrences for generating orthogonal Krylov subspace bases. SIAM Rev. 50, 485–503 (2008)

% block: reference-source-122 page:228 type:ordinary-paragraph
\item\label{bib:source-122}Liesen, J., Strakoš, Z.: Krylov Subspace Methods: Principles and Analysis. Oxford University Press, Oxford (2012)

% block: reference-source-123 page:228 type:ordinary-paragraph
\item\label{bib:source-123}Liu, J.W.H.: Modification of the minimum-degree algorithm by multiple elimination. ACM Trans. Math. Softw. 11, 141–153 (1985)

% block: reference-source-124 page:228 type:ordinary-paragraph
\item\label{bib:source-124}Liu, L., Sekiya, K., Ogino, M., Masui, K.: A COMINRES-QLP method for solving complex symmetric linear systems. Electr. Eng. Jpn. 214, e23325 (2021)

% block: reference-source-125 page:228 type:ordinary-paragraph
\item\label{bib:source-125}Lottes, J.: Towards Robust Algebraic Multigrid Methods for Nonsymmetric Problems. Springer, Cham (2017)

% block: reference-source-126 page:228 type:ordinary-paragraph
\item\label{bib:source-126}Luenberger, D.: Hyperbolic pairs in the method of conjugate gradients. SIAM J. Appl. Math. 6, 1263–1267 (1969)

% block: reference-source-127 page:228 type:ordinary-paragraph
\item\label{bib:source-127}Meijerink, J.A., van der Vorst, H.A.: An iterative solution method for linear systems of which the coefficient matrix is a symmetric M-matrix. Math. Comp. 31, 148–162 (1977)

% block: reference-source-128 page:228 type:ordinary-paragraph
\item\label{bib:source-128}Meng, G.-Y.: A practical asymptotical optimal SOR method. Appl. Math. Comput. 242, 707–715 (2014)

% block: reference-source-129 page:228 type:ordinary-paragraph
\item\label{bib:source-129}Meurant, G., Tebbens, J.D.: Krylov Methods for Nonsymmetric Linear Systems: From Theory to Computations, Springer Series in Computational Mathematics, vol. 57. Springer Nature, Switzerland (2020)

% block: reference-source-130 page:228 type:ordinary-paragraph
\item\label{bib:source-130}Miyatake, Y., Sogabe, T., Zhang, S.-L.: On the equivalence between SOR-type methods for linear systems and the discrete gradient methods for gradient systems. J. Comput. Appl. Math. 342, 58–69 (2018)

% block: reference-source-131 page:228 type:ordinary-paragraph
\item\label{bib:source-131}Miyatake, Y., Sogabe, T., Zhang, S.-L.: Adaptive SOR methods based on the Wolfe conditions. Numer. Algorithms 84, 117–132 (2020)

% block: reference-source-132 page:228 type:ordinary-paragraph
\item\label{bib:source-132}Mizuno, S., Moriizumi, Y., Usuda, T.S., Sogabe, T.: An initial guess of Newton’s method for the matrix square root based on a sphere constrained optimization problem. JSIAM Lett. 8, 17–20 (2016)

% block: reference-source-133 page:228 type:ordinary-paragraph
\item\label{bib:source-133}Moler, C., Van Loan, C.: Nineteen dubious ways to compute the exponential of a matrix, twenty-five years later. SIAM Rev. 45, 3–49 (2003)

% block: reference-source-134 page:228 type:ordinary-paragraph
\item\label{bib:source-134}Morgan, R.B.: A restarted GMRES method augmented with eigenvectors. SIAM J. Matrix Anal. Appl. 16, 1154–1171 (1995)

% block: reference-source-135 page:228 type:ordinary-paragraph
\item\label{bib:source-135}Morgan, R.B.: GMRES with deflated restarting. SIAM J. Sci. Comput. 24, 20–37 (2002)

% block: reference-source-136 page:228 type:ordinary-paragraph
\item\label{bib:source-136}Morikuni, K., Hayami, K.: Inner-iteration Krylov subspace methods for least squares problems. SIAM J. Matrix Anal. Appl. 34, 1–22 (2013)

% block: reference-source-137 page:228 type:ordinary-paragraph
\item\label{bib:source-137}Moriya, K., Nodera, T.: The DEFLATED-GMRES(m, k) method with switching the restart frequency dynamically. Numer. Linear Algebra Appl. 7, 569–584 (2000)

% block: reference-source-138 page:228 type:ordinary-paragraph
\item\label{bib:source-138}Nachtigal, N.M., Reddy, S.C., Trefethen, L.N.: How fast are nonsymmetric matrix iterations? SIAM J. Matrix Anal. Appl. 13, 778–795 (1992)

% block: reference-source-139 page:228 type:ordinary-paragraph
\item\label{bib:source-139}Noritake, S., Imakura, A., Yamamoto, Y., Zhang, S.-L.: A large-grained parallel solver for linear simultaneous ordinary differential equations based on matrix exponential and its evaluation (in Japanese). Trans. JSIAM 19, 293–312 (2009)

% block: reference-source-140 page:228 type:ordinary-paragraph
\item\label{bib:source-140}Ogasawara, M., Tadano, H., Sakurai, T., Itoh, S.: A Krylov subspace method for shifted linear systems and its application to eigenvalue problems (in Japanese). Trans. JSIAM 14, 193–205 (2004)

% block: reference-source-141 page:228 type:ordinary-paragraph
\item\label{bib:source-141}Ogino, M., Takei, A., Sugimoto, S., Yoshimura, S.: A numerical study of iterative substructuring method for finite element analysis of high frequency electromagnetic fields. Comput. Math. Appl. 72, 2020–2027 (2016)

% block: reference-source-142 page:228 type:ordinary-paragraph
\item\label{bib:source-142}O’Leary, D.P.: The block conjugate gradient algorithm and related methods. Linear Algebra Appl. 29, 293–322 (1980)

% block: reference-source-143 page:228 type:ordinary-paragraph
\item\label{bib:source-143}Paige, C.C., Saunders, M.A.: Solution of sparse indefinite systems of linear equations. SIAM J. Numer. Anal. 12, 617–629 (1975)

% block: reference-source-144 page:228 type:ordinary-paragraph
\item\label{bib:source-144}Paige, C.C., Saunders, M.A.: LSQR: An algorithm for sparse linear equations and sparse least squares. ACM Trans. Math. Softw. 8, 43–71 (1982)

% block: reference-source-145 page:229 type:ordinary-paragraph
\item\label{bib:source-145}Parlett, B.N., Taylor, D.R., Liu, Z.A.: A look-ahead Lanczos algorithm for unsymmetric matrices. Math. Comp. 44, 105–124 (1985)

% block: reference-source-146 page:229 type:ordinary-paragraph
\item\label{bib:source-146}Pearson, J.W., Pestana, J.: Preconditioners for Krylov subspace methods: An overview. GAMM-Mitteilungen 43, e202000015 (2020)

% block: reference-source-147 page:229 type:ordinary-paragraph
\item\label{bib:source-147}Pommerell, C.: Solution of Large Unsymmetric Systems of Linear Equations. Ph.D. thesis, ETH Zürich, Switzerland (1992)

% block: reference-source-148 page:229 type:ordinary-paragraph
\item\label{bib:source-148}Quinn, J.A., Sugiyama, M.: A least-squares approach to anomaly detection in static and sequential data. Pattern Recognit. Lett. 40, 36–40 (2014)

% block: reference-source-149 page:229 type:ordinary-paragraph
\item\label{bib:source-149}Saad, Y.: Krylov subspace methods for solving large unsymmetric linear systems. Math. Comp. 37, 105–126 (1981)

% block: reference-source-150 page:229 type:ordinary-paragraph
\item\label{bib:source-150}Saad, Y.: ILUT: A dual threshold incomplete LU factorization. Numer. Linear Algebra Appl. 1, 387–402 (1994)

% block: reference-source-151 page:229 type:ordinary-paragraph
\item\label{bib:source-151}Saad, Y.: Iterative Methods for Sparse Linear Systems, 2nd edn. SIAM, Philadelphia (2003)

% block: reference-source-152 page:229 type:ordinary-paragraph
\item\label{bib:source-152}Saad, Y., Schultz, M.H.: GMRES: A generalized minimal residual algorithm for solving nonsymmetric linear systems. SIAM J. Sci. Stat. Comput. 7, 856–869 (1986)

% block: reference-source-153 page:229 type:ordinary-paragraph
\item\label{bib:source-153}Saito, S., Tadano, H., Imakura, A.: Development of the block BiCGSTAB($\ell$) method for solving linear systems with multiple right hand sides. JSIAM Lett. 6, 65–68 (2014)

% block: reference-source-154 page:229 type:ordinary-paragraph
\item\label{bib:source-154}Sakurai, T., Matsuo, T., Katagiri, T. (eds.): Numerical Linear Algebra: Theory and HPC (in Japanese). Kyoritsu Shuppan, Tokyo (2018)

% block: reference-source-155 page:229 type:ordinary-paragraph
\item\label{bib:source-155}Sato, H., Aihara, K.: Riemannian Newton’s method on the Grassmann manifold exploiting the quotient structure (in Japanese). Trans. JSIAM 28, 205–241 (2018)

% block: reference-source-156 page:229 type:ordinary-paragraph
\item\label{bib:source-156}Sattler, K.D. (ed.): 21st Century Nanoscience—A Handbook: Nanophysics Source Book (1st ed.). CRS Press, Boca Raton (2020)

% block: reference-source-157 page:229 type:ordinary-paragraph
\item\label{bib:source-157}Simoncini, V.: A stabilized QMR version of block BiCG. SIAM J. Matrix Anal. Appl. 18, 419–434 (1997)

% block: reference-source-158 page:229 type:ordinary-paragraph
\item\label{bib:source-158}Simoncini, V.: Restarted full orthogonalization method for shifted linear systems. BIT Numer. Math. 43, 459–466 (2003)

% block: reference-source-159 page:229 type:ordinary-paragraph
\item\label{bib:source-159}Simoncini, V., Perotti, F.: On the numerical solution of $(\lambda^2 A+\lambda B+C)\boldsymbol{x}=\boldsymbol{b}$ and application to structural dynamics. SIAM J. Sci. Comput. 23, 1875–1897 (2002)

% block: reference-source-160 page:229 type:ordinary-paragraph
\item\label{bib:source-160}Simoncini, V., Szyld, D.B.: Recent computational developments in Krylov subspace methods for linear systems. Numer. Linear Algebra Appl. 14, 1–59 (2007)

% block: reference-source-161 page:229 type:ordinary-paragraph
\item\label{bib:source-161}Simoncini, V., Szyld, D.B.: Interpreting IDR as a Petrov-Galerkin method. SIAM J. Sci. Comput. 32, 1898–1912 (2010)

% block: reference-source-162 page:229 type:ordinary-paragraph
\item\label{bib:source-162}Sleijpen, G.L.G., Fokkema, D.R.: BiCGstab($\ell$) for linear equations involving unsymmetric matrices with complex spectrum. Electron. Trans. Numer. Anal. 1, 11–32 (1993)

% block: reference-source-163 page:229 type:ordinary-paragraph
\item\label{bib:source-163}Sleijpen, G.L.G., Sonneveld, P., van Gijzen, M.B.: Bi-CGSTAB as an induced dimension reduction method. Appl. Numer. Math. 60, 1100–1114 (2010)

% block: reference-source-164 page:229 type:ordinary-paragraph
\item\label{bib:source-164}Sleijpen, G.L.G., van Gijzen, M.B.: Exploiting BiCGstab($\ell$) strategies to induce dimension reduction. SIAM J. Sci. Comput. 32, 2687–2709 (2010)

% block: reference-source-165 page:229 type:ordinary-paragraph
\item\label{bib:source-165}Sogabe, T.: Extensions of the Conjugate Residual Method. Ph.D. thesis, The University of Tokyo, Tokyo, Japan (2006)

% block: reference-source-166 page:229 type:ordinary-paragraph
\item\label{bib:source-166}Sogabe, T., Hoshi, T., Zhang, S.-L., Fujiwara, T.: A numerical method for calculating the Green’s function arising from electronic structure theory. In: Kaneda, Y., Kawamura, H., Sasai, M. (eds.) Frontiers of Computational Science, pp. 189–195. Springer, Berlin, Heidelberg (2007)

% block: reference-source-167 page:229 type:ordinary-paragraph
\item\label{bib:source-167}Sogabe, T., Hoshi, T., Zhang, S.-L., Fujiwara, T.: On a weighted quasi-residual minimization strategy for solving complex symmetric shifted linear systems. Electron. Trans. Numer. Anal. 31, 126–140 (2008)

% block: reference-source-168 page:229 type:ordinary-paragraph
\item\label{bib:source-168}Sogabe, T., Hoshi, T., Zhang, S.-L., Fujiwara, T.: Solution of generalized shifted linear systems with complex symmetric matrices. J. Comput. Phys. 231, 5669–5684 (2012)

% block: reference-source-169 page:229 type:ordinary-paragraph
\item\label{bib:source-169}Sogabe, T., Sugihara, M., Zhang, S.-L.: An extension of the conjugate residual method to nonsymmetric linear systems. J. Comput. Appl. Math. 226, 103–113 (2009)

% block: reference-source-170 page:229 type:ordinary-paragraph
\item\label{bib:source-170}Sogabe, T., Yamamoto, Y.: Basic Mathematical Algorithm for Computational Science (in Japanese). Zhang, S.-L. (ed.), Kaneda, Y., (ed. in chief), Sasai, M (ed. in chief), Kyoritsu Shuppan, Tokyo (2019)

% block: reference-source-171 page:230 type:ordinary-paragraph
\item\label{bib:source-171}Sogabe, T., Zhang, S.-L.: A COCR method for solving complex symmetric linear systems. J. Comput. Appl. Math. 199, 297–303 (2007)

% block: reference-source-172 page:230 type:ordinary-paragraph
\item\label{bib:source-172}Sogabe, T., Zhang, S.-L.: An extension of the COCR method to solving shifted linear systems with complex symmetric matrices. E. Asian J. Appl. Math. 1, 97–107 (2011)

% block: reference-source-173 page:230 type:ordinary-paragraph
\item\label{bib:source-173}Sonneveld, P.: CGS, a fast Lanczos-type solver for nonsymmetric linear systems. SIAM J. Sci. Stat. Comput. 10, 36–52 (1989)

% block: reference-source-174 page:230 type:ordinary-paragraph
\item\label{bib:source-174}Sonneveld, P., van Gijzen, M.B.: IDR(s): A family of simple and fast algorithms for solving large nonsymmetric systems of linear equations. SIAM J. Sci. Comput. 31, 1035–1062 (2008)

% block: reference-source-175 page:230 type:ordinary-paragraph
\item\label{bib:source-175}Soodhalter, K.M., Szyld, D.B., Xue, F.: Krylov subspace recycling for sequences of shifted linear systems. Appl. Numer. Math. 81, 105–118 (2014)

% block: reference-source-176 page:230 type:ordinary-paragraph
\item\label{bib:source-176}Sosonkina, M., Watson, L.T., Kapania, R.K., Walker, H.F.: A new adaptive GMRES algorithm for achieving high accuracy. Numer. Linear Algebra Appl. 5, 275–297 (1998)

% block: reference-source-177 page:230 type:ordinary-paragraph
\item\label{bib:source-177}Stiefel, E.: Relaxationsmethoden bester strategie zur lösung linearer gleichungssysteme. Comment. Math. Helv. 29, 157–179 (1955)

% block: reference-source-178 page:230 type:ordinary-paragraph
\item\label{bib:source-178}Sugiyama, M.: Superfast-trainable multi-class probabilistic classifier by least-squares posterior fitting. IEICE Trans. Inform. Syst. E93-D, 2690–2701 (2010)

% block: reference-source-179 page:230 type:ordinary-paragraph
\item\label{bib:source-179}Suzuki, M.: Generalized Trotter’s formula and systematic approximants of exponential operators and inner derivations with applications to many-body problems. Communs Math. Phys. 51, 183–190 (1976)

% block: reference-source-180 page:230 type:ordinary-paragraph
\item\label{bib:source-180}Sylvester, J.J.: On the equation to the secular inequalities in the planetary theory. Philos. Mag. 16, 267–269 (1883)

% block: reference-source-181 page:230 type:ordinary-paragraph
\item\label{bib:source-181}Tadano, H.: Development of the Block BiCGGR$^2$ method for linear systems with multiple right-hand sides. Japan J. Indust. Appl. Math. 36, 563–577 (2019)

% block: reference-source-182 page:230 type:ordinary-paragraph
\item\label{bib:source-182}Tadano, H., Sakurai, T., Kuramashi, Y.: Block BiCGGR: A new block Krylov subspace method for computing high accuracy solutions. JSIAM Lett. 1, 44–47 (2009)

% block: reference-source-183 page:230 type:ordinary-paragraph
\item\label{bib:source-183}Takahasi, H., Mori, M.: Double exponential formulas for numerical integration. Publ. Res. Inst. Math. Sci. 9, 721–741 (1973)

% block: reference-source-184 page:230 type:ordinary-paragraph
\item\label{bib:source-184}Takayama, R., Hoshi, T., Sogabe, T., Zhang, S.-L., Fujiwara, T.: Linear algebraic calculation of the Green’s function for large-scale electronic structure theory. Phys. Rev. B 73, 165108 (2006)

% block: reference-source-185 page:230 type:ordinary-paragraph
\item\label{bib:source-185}Tanio, M., Sugihara, M.: GBi-CGSTAB(s, L): IDR(s) with higher-order stabilization polynomials. J. Comput. Appl. Math. 235, 765–784 (2010)

% block: reference-source-186 page:230 type:ordinary-paragraph
\item\label{bib:source-186}Tatsuoka, F., Sogabe, T., Miyatake, Y., Kemmochi, T., Zhang, S.-L.: Computing the matrix fractional power with the double exponential formula. Electron. Trans. Numer. Anal. 54, 558–580 (2021)

% block: reference-source-187 page:230 type:ordinary-paragraph
\item\label{bib:source-187}Tatsuoka, F., Sogabe, T., Miyatake, Y., Zhang, S.-L.: A cost-efficient variant of the incremental newton iteration for the matrix pth root. J. Math. Res. Appl. 37, 97–106 (2017)

% block: reference-source-188 page:230 type:ordinary-paragraph
\item\label{bib:source-188}Tatsuoka, F., Sogabe, T., Miyatake, Y., Zhang, S.-L.: Algorithms for the computation of the matrix logarithm based on the double exponential formula. J. Comput. Appl. Math. 373, 112396 (2020)

% block: reference-source-189 page:230 type:ordinary-paragraph
\item\label{bib:source-189}Taylor, D.: Analysis of the Look Ahead Lanczos Algorithm. Ph.D. thesis, University of California, Berkley (1982)

% block: reference-source-190 page:230 type:ordinary-paragraph
\item\label{bib:source-190}Teng, H., Fujiwara, T., Hoshi, T., Sogabe, T., Zhang, S.-L., Yamamoto, S.: Efficient and accurate linear algebraic methods for large-scale electronic structure calculations with nonorthogonal atomic orbitals. Phys. Rev. B 83, 165103 (2011)

% block: reference-source-191 page:230 type:ordinary-paragraph
\item\label{bib:source-191}Trefethen, L.N., Bau, D.: Numerical Linear Algebra: Twenty-Fifth. Anniversary SIAM, Philadelphia (2022)

% block: reference-source-192 page:230 type:ordinary-paragraph
\item\label{bib:source-192}Trefethen, L.N., Weideman, J.A.C.: The exponentially convergent trapezoidal rule. SIAM Rev. 56, 385–458 (2014)

% block: reference-source-193 page:230 type:ordinary-paragraph
\item\label{bib:source-193}Trottenberg, U., Oosterlee, C.W., Schüller, A.: Multigrid. Academic Press, NewYork (2001)

% block: reference-source-194 page:230 type:ordinary-paragraph
\item\label{bib:source-194}van der Sluis, A.: Condition numbers and equilibration of matrices. Numer. Math. 14, 14–23 (1969)

% block: reference-source-195 page:230 type:ordinary-paragraph
\item\label{bib:source-195}van der Vorst, H.A.: Bi-CGSTAB: A fast and smoothly converging variant of Bi-CG for the solution of nonsymmetric linear systems. SIAM J. Sci. Stat. Comput. 13, 631–644 (1992)

% block: reference-source-196 page:231 type:ordinary-paragraph
\item\label{bib:source-196}van der Vorst, H.A.: Iterative Krylov Methods for Large Linear Systems. Cambridge University Press, Cambridge (2003)

% block: reference-source-197 page:231 type:ordinary-paragraph
\item\label{bib:source-197}van der Vorst, H.A., Melissen, J.B.M.: A Petrov-Galerkin type method for solving $A\boldsymbol{x}=\boldsymbol{b}$, where $A$ is symmetric complex. IEEE Trans. Magn. 26, 706–708 (1990)

% block: reference-source-198 page:231 type:ordinary-paragraph
\item\label{bib:source-198}van Gijzen, M.B.: A polynomial preconditioner for the GMRES algorithm. J. Comput. Appl. Math. 59, 91–107 (1995)

% block: reference-source-199 page:231 type:ordinary-paragraph
\item\label{bib:source-199}van Gijzen, M.B., Sonneveld, P.: Algorithm 913: An elegant IDR(s) variant that efficiently exploits biorthogonality properties. ACM Trans. Math. Softw. 38, 5:1–5:19 (2011)

% block: reference-source-200 page:231 type:ordinary-paragraph
\item\label{bib:source-200}Varga, R.S.: Matrix Iterative Analysis, 2nd edn. Springer-Verlag, Berlin Heidelberg (2000)

% block: reference-source-201 page:231 type:ordinary-paragraph
\item\label{bib:source-201}Vinsome, P.K.W.: Orthomin, an iterative method for solving sparse sets of simultaneous linear equations. In: Fourth Symposium on Reservoir Simulation, pp. 149–159. Society of Petroleum Engineers of AIME (1976)

% block: reference-source-202 page:231 type:ordinary-paragraph
\item\label{bib:source-202}Vital, B.: Etude de Quelques Méthodes de Résolution de Problèmes Linéaires de Grande Taille sur Multiprocesseur. Ph.D. thesis, Univ. de Rennes I, Rennes, France (1990)

% block: reference-source-203 page:231 type:ordinary-paragraph
\item\label{bib:source-203}Weiss, R.: Parameter-Free Iterative Linear Solvers, 1st edn. Akademie Verlag, Berlin (1996)

% block: reference-source-204 page:231 type:ordinary-paragraph
\item\label{bib:source-204}Wesseling, P.: An Introduction to Multigrid Methods. Wiley, Chichester (1991)

% block: reference-source-205 page:231 type:ordinary-paragraph
\item\label{bib:source-205}Wesseling, P., Sonneveld, P.: Numerical experiments with a multiple grid and a preconditioned Lanczos type method. In: Rautmann, R. (ed.) Approximation Methods for Navier-Stokes Problems, pp. 543–562. Springer, Berlin, Heidelberg (1980)

% block: reference-source-206 page:231 type:ordinary-paragraph
\item\label{bib:source-206}Yin, J.-F., Hayami, K.: Preconditioned GMRES methods with incomplete Givens orthogonalization method for large sparse least-squares problems. J. Comput. Appl. Math. 226, 177–186 (2009)

% block: reference-source-207 page:231 type:ordinary-paragraph
\item\label{bib:source-207}Young, D.M.: The search for high level parallelism for the iterative solution of large sparse linear systems. In: Carey, G.F. (ed.) Parallel Supercomputing: Methods, Algorithms and Applications, pp. 89–105. Wiley, Chichester, UK (1989)

% block: reference-source-208 page:231 type:ordinary-paragraph
\item\label{bib:source-208}Young, D.M., Vona, B.R.: On the use of rational iterative methods for solving large sparse linear systems. Appl. Numer. Math. 10, 261–278 (1992)

% block: reference-source-209 page:231 type:ordinary-paragraph
\item\label{bib:source-209}Zemke, J.-P.M.: Variants of IDR with partial orthonormalization. Electron. Trans. Numer. Anal. 46, 245–272 (2017)

% block: reference-source-210 page:231 type:ordinary-paragraph
\item\label{bib:source-210}Zhang, J.H., Zhao, J.: A novel class of block methods based on the block $AA^{\top}$-Lanczos bi-orthogonalization process for matrix equations. Int. J. Comput. Math. 90, 341–359 (2013)

% block: reference-source-211 page:231 type:ordinary-paragraph
\item\label{bib:source-211}Zhang, L., Nodera, T.: A new adaptive restart for GMRES(m) method. ANZIAM J. 46, C409–C425 (2004)

% block: reference-source-212 page:231 type:ordinary-paragraph
\item\label{bib:source-212}Zhang, S.-L.: GPBi-CG: Generalized product-type methods based on Bi-CG for solving non-symmetric linear systems. SIAM J. Sci. Comput. 18, 537–551 (1997)

\end{enumerate}
